\section{Empirical Results}\label{res}

Our empirical analysis identifies statistically significant discontinuities at the threshold of closely contested auctions, suggesting potential anti-competitive dynamics in public procurement markets.

Graphical representations of these effects are illustrated in Figures \ref{fig:incumbent} and \ref{fig:lastbid}. Quantitative estimates supporting the graphical evidence are presented in Table \ref{tab:estimates}.
Figure \ref{fig:incumbent} shows the estimated probability that the winning firm was the incumbent—defined as the firm winning the immediately preceding auction for a similar good—as a function of the bid margin. A clear discontinuity at the threshold (bid margin equal to zero) indicates that firms narrowly winning auctions (negative bid margin) exhibit a higher probability of incumbency compared to narrowly losing firms (positive bid margin). This pattern is consistent with incumbent firms benefiting disproportionately, which may be reflective of strategic bidding practices or bid rotation schemes.

\begin{figure}[htbp!]
\centering
\caption{Probability of Incumbency around the Bid Margin}\label{fig:incumbent}
\includegraphics[width=0.8\textwidth]{Presentations/rd_plot_won_t_minus_1_overall_market.pdf}
\begin{minipage}{0.8\textwidth}
{\scriptsize \emph{Notes:} The figure plots binned averages of incumbency status against the bid margin, where negative values correspond to narrowly won auctions and positive values reflect narrowly lost auctions. A distinct discontinuity at the cutoff suggests potential incumbency advantages. \par}
\end{minipage}
\end{figure}

Similarly, Figure \ref{fig:lastbid} examines the probability that the winning firm submitted the last bid in the auction as a function of the bid margin. A comparable discontinuity at the cutoff suggests firms narrowly winning auctions have a higher likelihood of being the final bidder, potentially indicating strategic bid timing indicative of anti-competitive behavior, such as bid leakage or collusion.

\begin{figure}[htbp!]
\centering
\caption{Probability of Being the Last Bidder around the Bid Margin}\label{fig:lastbid}
\includegraphics[width=0.8\textwidth]{Presentations/rd_plot_last_bid_overall_market.pdf}
\begin{minipage}{0.8\textwidth}
{\scriptsize \emph{Notes:} The figure plots binned averages of the probability that the winning firm submitted the last bid against the bid margin, where negative values correspond to narrowly won auctions and positive values reflect narrowly lost auctions. A discontinuity at the threshold indicates potential strategic timing in bid submissions. \par}
\end{minipage}
\end{figure}

The regression discontinuity estimates corresponding to these graphical results, accounting also for year and market fixed effects, are reported in Table \ref{tab:estimates}. We find that being the incumbent is associated with a 1.82\% increase in the probability of winning an auction within the same market. Similarly, winning firms are  6.36\% more likely to submit the final bid. Both coefficients are statistically significant at the 1\% level. The results remain robust to the inclusion of market and year fixed effects and statistically support the graphical evidence, reinforcing the interpretation of non-competitive forces at play in closely contested procurement auctions.

\begin{table}[htbp]
\centering
\caption{Effect of Winning Auction on Key Outcomes}\label{tab:estimates}
\begin{tabular}{lcccc}
\toprule
 & \multicolumn{2}{c}{Incumbent} & \multicolumn{2}{c}{Last Bid} \\
\cmidrule(lr){2-3} \cmidrule(lr){4-5}
 & (1) & (2) & (3) & (4) \\
\midrule
Winning Auction & $0.0182^{***}$ & $0.0179^{***}$ & $0.0636^{***}$ & $0.0636^{***}$ \\
 & ($0.005$) & ($0.005$) & ($0.008$) & ($0.008$) \\
\\
Observations & $261370$ & $260985$ & $156418$ & $156418$ \\
Optimal Bandwidth & 0.02 & 0.02 & 0.01 & 0.01 \\
Year FE &  & \checkmark &  & \checkmark \\
Market FE &  & \checkmark &  & \checkmark \\
\bottomrule
\end{tabular}
\vspace{0.5em}
\begin{minipage}{0.9\textwidth}
\scriptsize{
\textbf{Notes:} Standard errors are reported in parentheses. *** indicates statistical significance at the 1\% level. The dependent variables are defined as follows: (1) \textit{Incumbent} is a dummy indicating whether the firm won the previous auction within the same market; (2) \textit{Last Bid} is a dummy variable indicating whether the winning firm submitted the final bid in the auction. \par}
\end{minipage}
\end{table}

\section{Additional Results - TBC}\label{res_ar}
\subsection{Heterogenety by market}
\subsection{Correlation with predicted corruption}
\subsection{Firms outcomes}
